\chapter{Multiple media types}
\label{ch:multimedia}

mybibli is not just for books. Four media
types\index{media types} are first-class citizens:

\begin{itemize}
  \item \textbf{Book} — the default. Single-volume novels, essays,
        reference works, children's picture books.
  \item \textbf{Comic / BD}\index{comic / BD} — single-issue comics
        or multi-position omnibus albums (e.g.\ a ``Tintin
        anthology'' that physically gathers volumes 1--5).
  \item \textbf{Audio} — CDs, vinyl records, audiobooks on
        physical media.
  \item \textbf{Film / series}\index{film} — DVDs, Blu-ray boxsets,
        TV-series collections.
\end{itemize}

The media type drives three things:

\begin{enumerate}
  \item \textbf{Which metadata provider is consulted first.} ISBNs
        route to BnF + Google Books + Open Library; films route to
        TMDb + OMDb; audio routes to MusicBrainz.
  \item \textbf{Which fields are exposed on the title page.}
        Audio titles surface a track list; films surface director,
        runtime, age rating.
  \item \textbf{The icon shown next to the title} in search results
        and the catalog list.
\end{enumerate}

\section{Books}

The default media type. Catalog as described in
section~\ref{sec:catalog-scan}. Multi-volume series (encyclopaedias,
trilogies, \dots) are modeled with a \emph{Series}\index{series}
entity in chapter~\ref{ch:configuration}.

\section{Comics and BD}

A single \textbf{omnibus}\index{omnibus} volume can contain several
positions in a series. mybibli supports this with the
\emph{multi-position}\index{multi-position} feature: at volume
creation time, you can set the volume's \emph{position range} (for
example, ``contains volumes 5--7''). Series-gap detection
(chapter~\ref{ch:usage}) understands these ranges and will not flag
positions 5, 6 or 7 as missing once the omnibus is on the shelf.

\section{Audio releases}

Audio titles use MusicBrainz\index{MusicBrainz} as the primary
metadata source. The barcode on the back of a CD (EAN-13) resolves
the same way as an ISBN.

\begin{itemize}
  \item \textbf{Tracks.} MusicBrainz returns the track list; mybibli
        renders it on the title page.
  \item \textbf{Format.} CD vs vinyl vs cassette is captured on the
        volume row (not on the title) — two identical pressings on
        different media are two volumes of the same title.
\end{itemize}

\section{Films and series}

Films and TV series use OMDb\index{OMDb} and TMDb\index{TMDb} as the
primary metadata sources. A series boxset is modeled as a
\emph{title} with a multi-position volume covering the disc range.

\begin{tip}
OMDb is free for low volumes (under 1000 requests per day); TMDb is
free with attribution. Either covers a household-sized
DVD / Blu-ray shelf. See chapter~\ref{ch:providers} for setup.
\end{tip}

\section{Picking the media type at scan time}

After you scan an unfamiliar barcode, mybibli guesses the media
type from the EAN prefix (the GS1 country prefix and the resolved
provider's response). The guess is shown on the catalog page; you
can override it before clicking \textbf{Add volume} if it is
wrong. Re-typing a title later is a few clicks
(\texttt{/title/<id>}) and the provider chain re-runs against the
new type if you trigger \emph{Re-fetch metadata}.
